\section{Опис процесу створення архітектури програмного комплексу}

Головною інженерною особливістю розробленого додатка є об'єднання двох різних 
за семантикою практичних задач в одному виконуваному файлі \texttt{Main.java} 
під керуванням віртуальної машини Java (JVM). Процес декомпозиції та побудови 
коду виконувався у кілька послідовних етапів з дотриманням принципу розділення 
відповідальності.

Для розуміння того, як буде виглядати взаємодія об'єктів у додатку, та облегшення
процесу створення (написання) коду, спочатку, за допомогою застосунку \texttt{Drawio},
побудуэмо UML-дыаграму.

Нижче на малюнку \ref{fig:AppDiagram_UML_lab5} наведено скріншот UML-діаграми для 
лабораторної роботи яка складається з двух окремих завдань, які об'єднані головним 
файлом \texttt{Main.java}.

\begin{figure}[H]
    \centering
    \includegraphics[width=0.95\textwidth]{screenshots/AppDiagram_UML_lab5.png}
    \caption{UML-діаграма класів лабораторної роботи №5}
    \label{fig:AppDiagram_UML_lab5}
\end{figure}

\subsection{Компонентний розбір архітектури системи}
Структура додатка чітко розмежована на два логічні пакети (packages), що запобігає 
змішуванню бізнес-контекстів і забезпечує чистоту проєктування:

\begin{enumerate}
    \item \textbf{Пакет \texttt{GeneralPracticalSection} (Загальний приклад):}
    \begin{itemize}
        \item \texttt{Vehicle} — базовий абстрактний клас, що задає інкапсульований 
        стан (\texttt{model}, \texttt{year}), відкриті геттери для доступу до них та
         абстрактний функціональний контракт у вигляді методу \texttt{move(): String}.
        \item \texttt{Car}, \texttt{Bicycle} — конкретні класи-спадкоємці, які 
        перевизначають метод \texttt{move()} для повернення специфічних рядків стану.
        \item \texttt{Chargeable} — функціональний інтерфейс-маркер поведінки із 
        методом \texttt{charge(): void}.
        \item \texttt{ElectricCar} — дочірній клас, який демонструє поєднання 
        вертикального успадкування (\texttt{extends Vehicle}) та горизонтального 
        розширення можливостей через імплементацію інтерфейсу (\texttt{implements Chargeable}).
    \end{itemize}
    
    \item \textbf{Пакет \texttt{IndividualAssignment} (Індивідуальний Варіант №7):}
    \begin{itemize}
        \item \texttt{ComputerDevice} — кореневий абстрактний клас ієрархії 
        обчислювальної техніки з полями \texttt{brand}, \texttt{model} та абстрактним 
        методом ініціалізації операційної системи \texttt{bootUp(): String}.
        \item \texttt{Connectable} — мережевий функціональний інтерфейс, який вводить 
        контракт підключення \texttt{connect(): void} та поліморфний метод опису 
        фізичного каналу зв'язку \texttt{getNetworkType(): String}.
        \item \texttt{Laptop}, \texttt{Desktop} — кінцеві конкретні класи сутностей. 
        Обидва класи одночасно розширюють базовий пристрій та імплементують мережевий 
        інтерфейс, надаючи унікальні геттери для власних прихованих полів (\texttt{batteryCapacity}
         та \texttt{powerSupply} відповідно).
    \end{itemize}
\end{enumerate}

\subsection{Аналіз типів зв'язків між елементами діаграми}
Для побудови коректної об'єктної моделі на схемі було використано три фундаментальні типи 
зв'язків стандарту UML:
\begin{itemize}
    \item \textbf{Успадкування / Генералізація (Generalization):} 
    Позначено суцільною лінією з незафарбованим трикутним покажчиком на кінці, спрямованим 
    від дочірніх класів до суперкласів (\texttt{Car} $\rightarrow$ \texttt{Vehicle}, 
    \texttt{Laptop} $\rightarrow$ \texttt{ComputerDevice}). Символізує відношення 
    \textit{«is-a»} (є цим об'єктом).
    \item \textbf{Реалізація / Імплементація (Realization):} Позначено пунктирною лінією з 
    незафарбованим трикутним покажчиком, спрямованим від конкретного класу до інтерфейсу 
    (\texttt{ElectricCar} $\dashrightarrow$ \texttt{Chargeable}, \texttt{Desktop}
     $\dashrightarrow$ \texttt{Connectable}). Визначає готовність класу виконувати 
     контракт поведінки інтерфейсу.
    \item \textbf{Залежність / Використання (Dependency):} Позначено простою стрілкою 
    від керуючого класу \texttt{Main} до масивів базових абстрактних типів. Вказує на 
    те, що зміни в сигнатурах методів \texttt{Vehicle} або \texttt{ComputerDevice} можуть 
    вплинути на логіку роботи точки входу, проте самі класи сутностей нічого не знають про
     існування класу \texttt{Main} (досягнуто низьку зв'язаність системи).
\end{itemize}

\subsection{Ізоляція бізнес-логіки сутностей від CLI-середовища}

Для забезпечення гнучкості архітектури було вирішено повністю відмовитися від 
використання прямих викликів \texttt{System.out.println()} всередині методів 
класів, що описують сутності. 

\begin{itemize}
    \item В ієрархії \textbf{Vehicle} абстрактний метод \texttt{move()} було 
    спроєктовано з сигнатурою, що повертає тип \texttt{String}. Таким чином, 
    класи \texttt{Car} та \texttt{Bicycle} лише формують та віддають текстовий 
    опис своєї дії, не зачіпаючи потоки виведення.
    \item В ієрархії \textbf{ComputerDevice} метод ініціалізації ОС \texttt{bootUp()} 
    аналогічно повертає \texttt{String}. Для доступу до внутрішнього стану 
    прихованих полів (\texttt{model}, \texttt{brand}, \texttt{batteryCapacity}, 
    \texttt{powerSupply}) було розроблено відкриті геттери, що дозволило зберегти 
    сувору інкапсуляцію даних.
\end{itemize}

\subsection{Проєктування та розширення інтерфейсних контрактів}

Інтерфейси в системі виступають як маркери додаткової поведінки, що застосовуються лише до певних вузлів ієрархії:
\begin{itemize}
    \item Інтерфейс \textbf{Chargeable} інтегрується виключно в 
    клас \texttt{ElectricCar}, вказуючи, що цей автомобіль потребує 
    зарядного пристрою.
    \item Мережевий інтерфейс \textbf{Connectable} за умовою 
    варіанта реалізують обидва класи комп'ютерної техніки (\texttt{Laptop} 
    та \texttt{Desktop}). Однак фізична природа підключення у них різна. 
    Щоб точка входу могла поліморфно дізнатися тип мережі, в контракт 
    \texttt{Connectable} було введено метод \texttt{getNetworkType(): String}, 
    який повертає \texttt{"Wi-Fi"} для портативних пристроїв та \texttt{"Ethernet"} 
    для стаціонарних ПК.
\end{itemize}

\subsection{Синхронізація обробки даних у \texttt{Main.java}}
Клас \texttt{Main.java} виступає в ролі центрального контролера та єдиного середовища 
представлення даних. Виконання двох завдань реалізовано через послідовні незалежні 
цикли обходу поліморфних масивів:
\begin{enumerate}
    \item \textbf{Блок загального завдання:} 
    Створюється масив посилань \texttt{Vehicle[]}, куди поміщаються об'єкти автомобіля, 
    велосипеда та електрокара. Під час ітерації додаток динамічно перевіряє підтримку 
    інтерфейсу \texttt{Chargeable} за допомогою оператора \texttt{instanceof} і виводить 
    статус зарядки.
    \item \textbf{Блок індивідуального варіанта:} 
    Створюється масив посилань \texttt{ComputerDevice[]}, який містить реальні екземпляри 
    ноутбуків та десктопів. Головний клас зчитує геттери бренду та моделі, динамічно 
    перевіряє наявність \texttt{Connectable} для отримання типу мережі, а також застосовує 
    безпечне низхідне приведення типів для отримання ємності батареї або потужності БЖ.
\end{enumerate}

Окрім циклів, у \texttt{Main.java} продемонстровано третій важливий аспект ООП: створення 
об'єкта \texttt{HP EliteBook} через інтерфейсну змінну типу \\
\texttt{Connectable officeLaptop = new Laptop(...)}. Це наочно ілюструє приховування 
деталей реалізації, коли програміст бачить лише методи підключення до мережі, але об'єкт 
у кучі залишається повноцінним ноутбуком, який згодом безпечно приводиться до базового 
класу для відображення в загальній структурі.